% EcoMD @ NeurIPS Workshop on Machine Learning and the Physical Sciences (ML4PS; non-archival, 4pp).
% Build: latexmk -pdf main.tex   (compiles with a NeurIPS-approximating fallback if neurips_2024.sty absent).
\documentclass{article}

\IfFileExists{neurips_2024.sty}{%
  \usepackage{neurips_2024}            % submission; add [final] for camera-ready
}{%
  \usepackage[letterpaper,margin=1in]{geometry}
  \usepackage{times}
  \usepackage{titlesec}
  \titleformat{\section}{\normalfont\large\bfseries}{\thesection}{0.6em}{}
  \titleformat{\subsection}{\normalfont\normalsize\bfseries}{\thesubsection}{0.6em}{}
  \setlength{\parskip}{2pt}
}

\usepackage[utf8]{inputenc}
\usepackage[T1]{fontenc}
\usepackage{amsmath,amssymb}
\usepackage{graphicx}
\usepackage{booktabs}
\usepackage{xcolor}
\usepackage{microtype}
\usepackage[numbers,sort&compress]{natbib}
\usepackage{hyperref}
\hypersetup{colorlinks=true,linkcolor=blue,citecolor=blue,urlcolor=blue}

\newcommand{\figbox}[2]{%
  \IfFileExists{#1}{\includegraphics[width=\linewidth]{#1}}{%
  \fbox{\parbox[c][3.0cm][c]{0.96\linewidth}{\centering\small\itshape #2}}}}

\title{Heavy Tails as a Non-Equilibrium Transient\\ in a Differentiable Particle Market Simulator}
\author{Anonymous Submission}
\date{}

\begin{document}
\maketitle

\begin{abstract}
Fat-tailed returns are usually treated as a \emph{stationary} property---the inverse-cubic tail
\citep{gabaix2003,plerou1999}. Using \textbf{EcoMD}, a differentiable Langevin particle simulator of a
market, we find the opposite. The \textbf{steady state is light-tailed} (Hill index $\alpha\!\approx\!4.7$),
and heavy, cube-law-and-beyond tails ($\alpha\!\to\!0.5$) appear \emph{only} when the system is driven
out of equilibrium---at initialization or by a controlled shock---and \textbf{relax on a finite
timescale} ($\tau\!\approx\!240$ steps). This makes warm-up-inclusive rollout scoring \emph{misread the
equilibration transient as a stationary law}. The transient requires \emph{coherent displacement of
agent latent states}: it is inert to an exogenous price shock (in both the return tail and order-flow
imbalance), and it leaves a sharp non-equilibrium order-flow signature (flow memory $0.02\!\to\!0.99$
at the shock, then relaxing). On real one-minute crypto crashes, by contrast, the return tail is
\emph{stationary} across calm and crash (a pre-registered null test over five episodes). The
simulator's heavy tail is thus a non-equilibrium phenomenon \emph{of the model}---localizing a missing
stationary heavy-tail mechanism in this class of simulators---rather than a stationary market law.
\end{abstract}

\section{Introduction}
The heavy tail of financial returns---the inverse-cubic law, $P(|r|\!>\!x)\!\sim\!x^{-\alpha}$ with
$\alpha\!\approx\!3$---is one of the most robust ``stylized facts'' in quantitative finance
\citep{cont2001,gabaix2003,plerou1999}, and is almost universally treated as a \emph{stationary}
property of the return distribution. Mechanistic, particle-based simulators offer a way to ask a
sharper, physics-style question: in a driven many-body system whose macroscopic observable is the
price, is the heavy tail a stationary equilibrium property, or a \emph{non-equilibrium transient}---and
does the standard way of measuring it tell the two apart?

We study this in \textbf{EcoMD}, a differentiable simulator in which agents are particles evolving
under overdamped Langevin dynamics with learned interaction potentials, and the price forms from
aggregate excess demand. Langevin and differentiable agent-based market models are not new
\citep{bouchaud1998,dyer2024,chopra2023}; our contribution is a non-equilibrium characterization and a
measurement correction:
\begin{itemize}\setlength{\itemsep}{1pt}
  \item Heavy tails are a \textbf{driven non-equilibrium transient} (light steady state; shock drives
    $\alpha\!\to\!0.5$; relaxation $\tau\!\approx\!240$), with a sigmoid dose-response over five assets
    (\S\ref{sec:transient}).
  \item A \textbf{non-equilibrium measurement pitfall}: warm-up-inclusive scoring confounds the
    equilibration transient with a stationary tail (\S\ref{sec:measurement}).
  \item \textbf{Mechanistic specificity and an order-flow signature}: only coherent latent displacement
    drives the transient (price shocks are inert), with a clean non-equilibrium order-flow signature
    (\S\ref{sec:mechanism}); and real return tails are \emph{stationary} (\S\ref{sec:boundary}).
\end{itemize}

\section{EcoMD: a Langevin particle market}\label{sec:ecomd}
Agents $s_i\!\in\!\mathbb{R}^d$ follow overdamped Langevin dynamics
$\dot s_i = -\nabla_{s_i}U(\{s\}) + \sqrt{2T}\,\xi_i$, with a learned equivariant interaction potential
$U$ (a message-passing form in the spirit of MACE \citep{batatia2022}) and temperature $T$. The
log-price increment is $\Delta p_t = \beta\,\mathrm{ED}_t + \sigma\eta_t$, where the (pre-impact)
excess demand $\mathrm{ED}_t=\kappa\sum_i\Delta s_{i,0}$ is the net signed agent flow; its Hill tail
index $\alpha_{\mathrm{ED}}$ \citep{hill1975} is the order-flow tail and, through the impact map,
controls the return tail.
The model is trained to reproduce a panel of stylized facts. Crucially, it is fully differentiable and
supports \emph{controlled interventions}---a scheduled shock that drives the system out of its steady
state---which we use as the non-equilibrium probe. We stress that because EcoMD \emph{is} a driven
overdamped Langevin many-body system, the terms ``steady state,'' ``relaxation time,'' and ``driven
transient'' below are precise statements about this dynamical system, not analogies for real markets;
the real-market question is treated separately and honestly in \S\ref{sec:boundary}.

\section{Heavy tails are a driven transient}\label{sec:transient}
In the unperturbed steady state the order-flow tail is \emph{light}: $\alpha_{\mathrm{ED}}\!\approx\!4.7$,
flat across the rollout (Fig.~\ref{fig:transient}, left, grey). We then apply a controlled
\emph{coherent displacement} (a \texttt{state\_kick}: a fraction of agents are displaced together) at
$t\!=\!3000$. The tail \textbf{craters to $\alpha_{\mathrm{ED}}\!\approx\!0.5$}---heavier than the cube
law---and then \textbf{relaxes back to $\approx\!4.7$ within $\sim\!10^3$ steps}. The depth matches the
$t\!=\!0$ equilibration template, i.e. the initial burn-in \emph{is} this same transient. An
exponential fit of the recovery limb gives a relaxation time $\tau\!\approx\!240$ steps for the
saturated regime (vs.\ $\tau\!\approx\!20$ for the faster cold-start). The effect generalizes across
five assets (equities, gold, crypto, FX) with a \textbf{sigmoid dose-response} in the shock magnitude
(onset $\sim\!0.1$--$0.2\sigma$, saturating at a floor $\alpha\!\approx\!0.5$; Fig.~\ref{fig:transient},
right). The heavy tail is therefore a property of the \emph{driven}, not the steady, state. Little of
this is imposed by the perturbation: the steady state being \emph{light} (the opposite of real
markets, and not built in), the finite relaxation time, the graded sigmoid dose-response, and the
mechanism-specificity of \S\ref{sec:mechanism} are all emergent properties of the trained dynamics.

\begin{figure}[t]
\centering
\figbox{figures/fig1_transient.pdf}{Fig.~1 placeholder --- (left) order-flow tail $\alpha_{\mathrm{ED}}(t)$:
control stays light ($\approx\!4.7$); a coherent shock at $t\!=\!3000$ craters it to $\approx\!0.5$
and it relaxes back ($\tau\!\approx\!240$). (right) sigmoid dose-response: post-shock minimum
$\alpha_{\mathrm{ED}}$ vs.\ shock magnitude.}
\caption{Heavy tails are a driven non-equilibrium transient: light steady state, shock-driven heavy
tail, finite-time relaxation, with a sigmoid dose-response.}
\label{fig:transient}
\end{figure}

\section{A non-equilibrium measurement pitfall}\label{sec:measurement}
Because a fresh rollout is out of equilibrium, it carries this heavy transient during its initial
relaxation. Standard stylized-fact scoring computes the Hill index on the \emph{full} rollout, with no
warm-up discard---so it reads the equilibration transient as a stationary heavy tail. Discarding the
first $\sim\!50$ steps moves the estimate sharply toward the light steady state: the baseline Hill
index rises $3.87\!\to\!6.55$, and a concave-impact variant previously reported to ``match'' the cube
law rises $5.05\!\to\!9.26$ (Fig.~\ref{fig:mech}a). The low cross-seed variance makes the artifact look
like a robust law. The correction is to score after a warm-up discard and report a warm-up-sensitivity
curve; under it, the earlier ``cube-law reproduction'' is withdrawn. This is a stationarity-hygiene
point for measuring tails in any driven simulator.

\section{Mechanism and order-flow signature}\label{sec:mechanism}
What drives the transient? Only \emph{coherent latent displacement}. An exogenous price shock (injecting
a return into the price---a news-shock analogue) is \textbf{inert}: across magnitudes up to $12\sigma$
the order-flow tail stays light ($\alpha_{\mathrm{ED}}\!\approx\!4.6$), the return tail only nudges
($\approx\!7.7\!\to\!6.0$, still far from heavy), and the order-flow imbalance is statistically
identical to the unshocked control. Under the coherent shock, order flow shows a clean non-equilibrium
signature: the lag-1 autocorrelation of order-flow imbalance jumps from $0.02$ to $\mathbf{0.99}$ at
the shock---flow becomes transiently coherent and persistent---and relaxes back, with the imbalance
magnitude and its extreme-coordination saturation bursting and decaying in step
(Fig.~\ref{fig:mech}b). This relaxing coherence burst is the order-flow fingerprint of the transient.
A driven, relaxing, time-asymmetric signal is exactly the regime where stochastic thermodynamics
\citep{seifert2012} predicts non-zero entropy production; quantifying the time-asymmetry and entropy
production of this signal---and comparing to fluctuation-theorem studies of market cascades
\citep{maskawa2025}---is a concrete next step that we do not claim to have established here.

\begin{figure}[t]
\centering
\figbox{figures/fig2_mech_boundary.pdf}{Fig.~2 placeholder --- (a) Hill index vs.\ warm-up discard:
warm-up-inclusive scoring reads a ``cube-law'' tail that vanishes once the transient is dropped.
(b) order-flow-imbalance memory bursts ($0.02\!\to\!0.99$) and relaxes under coherent shock, but is
flat under an exogenous price gap. (c) real-crash $\Delta\alpha$ against the calm null --- real return
tails do not heavy-up at crashes.}
\caption{(a) the measurement pitfall; (b) the non-equilibrium order-flow signature and its
mechanism-specificity; (c) the real-data stationarity boundary.}
\label{fig:mech}
\end{figure}

\section{The real-data boundary}\label{sec:boundary}
Is the same transient present in \emph{real} markets? We test it directly. On five real one-minute
crypto crash episodes (COVID-2020, the May-2021 selloff, the June-2022 deleveraging, Terra/Luna, FTX),
we pre-registered the statistic $\Delta\alpha = \alpha(\text{crash})-\alpha(\text{pre})$ on
volatility-standardized returns and compared it to a null distribution from a long calm window. The
driven-transient hypothesis predicts $\Delta\alpha\!\ll\!0$ (a heavier tail at the crash). Instead the
pooled effect is $z\!=\!+1.03$ (slightly \emph{lighter}); only one of five episodes is a significant
heavier-tail hit, within the false-positive rate for five tests (Fig.~\ref{fig:mech}c). Real return
tails are \textbf{stationary}---approximately cube-law in calm and crash alike, consistent with the
inverse-cubic universality \citep{gabaix2003,toth2018}. The simulator's transient heavy tail is
therefore a non-equilibrium property \emph{of the model}: because EcoMD has no stationary heavy-tail
source---a structural property, not a training shortfall---it can produce heavy tails only
transiently, and the light steady state is itself the finding. This localizes the missing
ingredient---a stationary heavy-tail mechanism (e.g.\ a heterogeneous, heavy order-flow source)---for
this class of simulators.

\section{Discussion}
We have shown, in a differentiable Langevin market simulator, that heavy tails are a \emph{driven
non-equilibrium transient} rather than a stationary law: a light steady state, a shock-driven heavy
tail that relaxes with a finite $\tau$, a sigmoid dose-response, and a relaxing order-flow coherence
signature---together with the measurement caveat that warm-up-inclusive scoring confounds this
transient with stationarity. The honest boundary is that \emph{real} return tails are stationary, so
the phenomenon is a property of this model class and pinpoints what it lacks. The relaxing order-flow
signature is a concrete, falsifiable prediction for real crash order-flow, a natural next test on
limit-order-book data. \textbf{Limitations.} The dynamical results are from a single simulator; we
conjecture the warm-up measurement pitfall affects any driven market simulator initialized off its
steady state, but verifying this across model families is future work. The real-data test is limited
to free intraday (crypto) data; the decisive follow-up is an order-flow-level test on paid
limit-order-book data, where the coherence signature above is the quantity to look for.

{\small
\bibliographystyle{plainnat}
\bibliography{references}
}

\end{document}
